\chapter{SYSTEM ANALYSIS AND DESIGN}

\section{System Requirements}

\subsection{Functional Requirements}
To address the challenges identified in the data collection workflow, the CTU-SignBridge platform must fulfill the following core functional requirements:
\begin{itemize}
    \item \textbf{Taxonomy Management:} Users must be able to create, read, update, and soft-delete vocabulary items within a highly structured, dialect-aware taxonomy.
    \item \textbf{Data Collection and Processing:} The system must allow users to record sign language videos via webcam, extract MediaPipe landmarks in real-time within the browser, and upload the compressed \texttt{.npz} arrays to the server.
    \item \textbf{Multi-tenant Workspace Management:} Administrators must be able to create isolated Workspaces for different research teams. Users must only see and interact with data belonging to their assigned Workspace.
    \item \textbf{Authentication and Authorization:} The system must support secure user login (SSO) and enforce strict Role-Based Access Control (RBAC) to ensure data integrity and prevent unauthorized modifications.
    \item \textbf{Asynchronous Synchronization:} Uploaded data must be reliably synchronized to cloud storage backends (e.g., Google Drive, MinIO) without blocking the user interface, complete with automatic retries on failure.
\end{itemize}

\subsection{Non-Functional Requirements}
The platform must also satisfy stringent non-functional requirements to ensure usability and scalability:
\begin{itemize}
    \item \textbf{Performance:} The API must achieve a P95 response time of $\leq 100$\,ms under a concurrent load of 500 virtual users.
    \item \textbf{Storage Efficiency:} The client-side extracted landmark data must reduce storage and bandwidth consumption by at least 90\% compared to raw MP4 video uploads.
    \item \textbf{Reliability:} Background synchronization tasks must achieve a $\geq 99$\,\% retry success rate in the event of simulated network failures.
\end{itemize}

\section{Overall System Architecture}

\subsection{Five-Layer Clean Architecture}
The backend is structured into five strict layers following the principle of separation of concerns \cite{dijkstra1982role}. Figure~\ref{fig:layered_arch} illustrates the data-flow direction between layers.

\insertfigure{0.85}{Pictures/layered_arch}{Five-layer Clean Architecture of CTU-SignBridge backend. Data flows inward from the HTTP boundary to the persistence layer and back.}{layered_arch}

\begin{enumerate}
    \item \textbf{Routers (HTTP Gateway):} Accept HTTP requests, validate JWT tokens (issued by Keycloak) via FastAPI dependency injection, and delegate to the Service layer. Contain \textit{no} business logic.
    \item \textbf{Schemas (Input/Output Validation):} Pydantic v2 models enforce type-safe contracts. Invalid payloads are rejected with HTTP 422 before reaching any business logic.
    \item \textbf{Services (Business Logic):} Orchestrate domain rules -- e.g., the 3-step Handshake commit, SHA-256 duplicate detection, and soft-delete enforcement. Casbin enforcers are invoked here to verify permissions.
    \item \textbf{Repositories (Data Access):} The sole layer authorized to issue SQL against PostgreSQL or call MinIO/Google Drive APIs.
    \item \textbf{Workers (Asynchronous Tasks):} Celery tasks handle cloud sync and heavy I/O operations, ensuring API latency is unaffected.
\end{enumerate}

\subsection{System Components and Data Flow}
The frontend (React/Vite) captures webcam streams and utilizes MediaPipe WebAssembly (Wasm) to compute landmarks. The resulting data is dispatched via HTTPS to the FastAPI backend. The backend persists metadata to PostgreSQL, offloads binary objects to MinIO, and delegates synchronization tasks to Celery via Redis. Authentication flows are entirely handled by redirecting the user to Keycloak, which returns a secure JWT upon successful login.

\section{Database Design}

\subsection{Entity-Relationship Diagram (ERD)}
The database is partitioned into distinct logical domains. Figure~\ref{fig:erd} illustrates the simplified ERD.

\insertfigure{0.95}{Pictures/erd}{Simplified ERD of CTU-SignBridge. Dotted arrowheads indicate nullable foreign keys.}{erd}

Key domains include:
\begin{itemize}
    \item \textbf{Auth \& Audit:} ROLES, USERS, WORKSPACES, WORKSPACE\_MEMBERS, PROJECTS.
    \item \textbf{Taxonomy (Star Schema):} LANGUAGES, DIALECTS, CATEGORIES, SIGN\_FEATURES, CLASSES.
    \item \textbf{Collection \& Media:} COLLECTION\_SESSIONS, SAMPLES, SAMPLE\_MEDIA, SAMPLE\_SYNC\_STATUS.
\end{itemize}

Workspace-level isolation (Multi-tenancy) is enforced at the data model layer: the PROJECTS table includes a non-nullable \texttt{workspace\_id} foreign key. All downstream tables (COLLECTION\_SESSIONS, SAMPLES, DATASETS) inherit this workspace scope through their foreign-key chain. 

\subsection{Star Schema Implementation}
To facilitate high-performance queries across the sign language vocabulary, the taxonomy is modeled as a Star Schema. The dimension tables (DIALECTS, CATEGORIES, SIGN\_FEATURES) are linked directly to the central CLASSES fact table via foreign keys (Table~\ref{tab:star_schema_fks}).

\begin{table}[H]
\centering
\caption{\textit{Foreign keys on the CLASSES table (Star Schema centre)}}
\label{tab:star_schema_fks}
\small
\begin{tabular}{p{3.2cm} p{3.0cm} p{1.8cm} p{4.5cm}}
\toprule
\textbf{Column} & \textbf{References} & \textbf{Nullable} & \textbf{Purpose} \\
\midrule
\texttt{dialect\_code} & DIALECTS & No  & Regional variant \\
\texttt{category\_id}  & CATEGORIES & Yes & Thematic grouping \\
\texttt{feature\_id}   & SIGN\_FEATURES & Yes & Gesture constraints \\
\bottomrule
\end{tabular}
\end{table}

\section{Security and Access Control Design}

\subsection{Integrating Keycloak for Authentication}
Authentication in CTU-SignBridge is offloaded to Keycloak using the OAuth 2.0 / OpenID Connect (OIDC) protocol. When a user attempts to access the platform, they are redirected to the Keycloak login portal. Upon providing valid credentials, Keycloak issues a signed JSON Web Token (JWT). The React frontend attaches this JWT as a Bearer token in the Authorization header of all subsequent API requests. The FastAPI backend validates the JWT's signature against Keycloak's public keys to confirm the user's identity securely without storing passwords locally.

\subsection{Access Enforcement using Casbin}
While Keycloak verifies identity, Casbin governs authorization. Casbin is integrated into the FastAPI Service layer to enforce Role-Based Access Control (RBAC) across tenants (Workspaces). 
The policy definitions in Casbin are mapped to Workspace memberships. For example, a user might hold the \texttt{Admin} role in \texttt{Workspace A} but only the \texttt{Viewer} role in \texttt{Workspace B}. Before executing any data-modifying operation, the Service layer invokes the Casbin Enforcer:
\texttt{enforcer.enforce(user\_id, workspace\_id, resource, action)}. 
If the evaluation returns false, the API immediately aborts the transaction and returns an HTTP 403 Forbidden error, guaranteeing strict multi-tenant isolation.

\section{Data Ingestion and Synchronization Workflow}

\subsection{Three-Step Handshake Protocol}
To prevent orphaned data and handle network instability during collection, a 3-step Handshake protocol is utilized (Figure~\ref{fig:live_capture}).

\insertfigure{0.75}{Pictures/live_capture}{Sequence diagram of the 3-step Handshake Commit protocol for Live Capture sessions.}{live_capture}

\begin{enumerate}
    \item \textbf{Create Session:} The client calls \texttt{POST /sessions} to create a COLLECTION\_SESSION with \texttt{status = in\_progress}.
    \item \textbf{Upload Data:} The client streams MediaPipe \texttt{.npz} objects directly to the MinIO object store via secure, temporary pre-signed URLs.
    \item \textbf{Commit Session:} The client calls \texttt{PATCH /sessions/\{id\}/commit}. This atomic database transaction transitions the session status to \texttt{completed} and bulk-inserts all corresponding SAMPLE metadata rows.
\end{enumerate}

\subsection{Asynchronous Cloud Sync (Celery, MinIO, Redis)}
Once a session is committed, the data must be securely backed up. The FastAPI server immediately responds to the user, ensuring a snappy UI experience, while a Celery task is enqueued in Redis. 
The Celery worker picks up the task from Redis, retrieves the binary \texttt{.npz} data from MinIO, and uploads it to an external cloud provider (e.g., Google Drive) via its API. If the upload fails due to network timeouts, the task is automatically re-queued in Redis with exponential back-off (e.g., retrying after 2s, 4s, 8s), ensuring eventual consistency and zero data loss.
